\documentclass[12pt]{article}
\usepackage[margin=1in]{geometry}
\usepackage[utf8]{inputenc}
\usepackage{latexsym, amsfonts, enumitem, graphicx, environ,enumitem, fancyhdr, color, amssymb, amsthm, amsmath, mathtools, tikz, nicematrix}
\usepackage{physics}
\usetikzlibrary{patterns}
\usetikzlibrary{decorations.pathreplacing}
\pagestyle{fancy}
\setlength{\headheight}{65pt}
\newenvironment{problem}[2][Problem]{\begin{trivlist}
    \item[\hskip \labelsep {\bfseries #1}\hskip \labelsep {\bfseries #2.}]}{\end{trivlist}}
\DeclareMathOperator{\Ima}{Im}

\usepackage{hyperref}
\usepackage[capitalize]{cleveref}
\usepackage{thmtools}

\newtheorem{thm}{Theorem}[section]
\newtheorem{lem}[thm]{Lemma}
\newtheorem{cl}[thm]{Claim}
\newtheorem{prop}[thm]{Proposition}
\newtheorem{cor}[thm]{Corollary}
\newtheorem{conj}[thm]{Conjecture}
\newtheorem{obstruction}[thm]{Obstruction}

\usepackage{mdframed}

\theoremstyle{definition}
\newtheorem{definition}[thm]{Definition}
\newtheorem{remark}[thm]{Remark}
\newtheorem{question}{Question}
\newtheorem{example}[thm]{Example}
\newtheorem{exercise}[thm]{Exercise}
\newtheorem{properties}[thm]{Properties}
\newtheorem*{answer}{Answer}

\usepackage{tcolorbox}
\tcbuselibrary{theorems,skins,breakable}

\tcolorboxenvironment{example}{
enhanced jigsaw,colframe=cyan,interior hidden,
breakable,%before skip=10pt,after skip=10pt
}

\tcolorboxenvironment{thm}{
enhanced jigsaw,colframe=red,interior hidden,
breakable,%before skip=10pt,after skip=10pt
}

\tcolorboxenvironment{prop}{
enhanced jigsaw,colframe=red,interior hidden,
breakable,%before skip=10pt,after skip=10pt
}

\tcolorboxenvironment{properties}{
enhanced jigsaw,colframe=red,interior hidden,
breakable,%before skip=10pt,after skip=10pt
}

\tcolorboxenvironment{lem}{
enhanced jigsaw,colframe=red,interior hidden,
breakable,%before skip=10pt,after skip=10pt
}
\tcolorboxenvironment{cor}{
enhanced jigsaw,colframe=red,interior hidden,
breakable,%before skip=10pt,after skip=10pt
}
\tcolorboxenvironment{obstruction}{
enhanced jigsaw,colframe=red,interior hidden,
breakable,%before skip=10pt,after skip=10pt
}

\tcolorboxenvironment{proof}{
enhanced jigsaw, frame hidden,%interior hidden,
breakable,%before skip=10pt,after skip=10pt
}

\tcolorboxenvironment{answer}{
enhanced jigsaw, frame hidden,%interior hidden,
breakable,%before skip=10pt,after skip=10pt
}

\renewcommand\qedsymbol{\( \blacksquare \)}

\lhead{Avi Chetlin \\ Assignment 03 \\ 14 September 2026}
\rhead{Dr. Stanislaw Szarek \\ MATH 423 Measure Theory \\ Fall 2026}

\begin{document}
All numbered problems from Folland, \textit{Real Analysis} (2e).

\begin{problem}{1.3.14}
  If \(\mu\) is a semifinite measure and \(\mu(E) = \infty\), for any \( C > 0 \) there exists \( F \subset E \) with \( C < \mu(F) < \infty \).
\end{problem}
\begin{proof}
  Consider the collection \( M = \{ \mu(F) \mid F \subset E, \mu(F) < \infty \} \), and define \( m \coloneq \sup M  \).
  Assume, indirectly, that there exists some \( C > 0 \) such that \( m \leq C < \infty \).
  By definition, \( m \) is a least upper bound.
  That is, for all \( n \in \mathbb{N} \), there exists some \(\mu\)-finite subset \( F_n \) of \( E \) which satisfies
  \begin{equation}
  \label{eq:5}
  m - \frac{1}{n} < \mu(F_n) < m.
  \end{equation}

  Define \( H_n = \bigcup_1^n F_n \), so that \( H_1 \subset H_2 \subset \dots \subset H_n \), and thus by monotonicity,
  \begin{equation}
  \label{eq:6}
  m - \frac{1}{n} < \mu(F_n) \leq \mu(H_n) \leq m.
  \end{equation}

  Finally, define \( F = \bigcup_1^{\infty} H_n \) as the union of the increasing sequence.
  By continuity from above, \( \mu(F) = \lim_{n\rightarrow \infty} \mu(H_n) \), which by \cref{eq:6} implies \(\mu(F) = m\).

  Now, consider the set \( E \setminus F \), which must be \(\mu\)-infinite since \( F \) is finite.
  Hence, it contains some \( A \subset E \setminus F \) with \(0 < \mu(A) < \infty\).
  Moreover, it is evident that \( \mu(A \cup F) \in M \).
  However, since \( A \) is disjoint with \( F \) by construction, \(\mu(A\cup F) = \mu(A) + \mu(F) > \mu(F) = m\), a contradiction.

\end{proof}

\begin{problem}{1.5.29}
  Let \( E \) be a Lebesgue measureable set.
  \begin{enumerate}[label=(\alph*)]
    \item If \( E \subset N \) where \( N \) is the nonmeasureable set described in \S 1.1, then \( m(E) = 0 \).
          \begin{proof}
          \end{proof}
    \item If \( m(E) > 0 \), then \( E \) contains a nonmeasureable set.
          (It suffices to assume \( E \subset [0,1] \). In the notation of \S 1.1, \( E = \bigcup_{r \in R} E \cap N_r \).)
  \end{enumerate}
\end{problem}

\begin{problem}{from class}
  Complete the details of the proof of the following fact.
  Let \( [a,b] \subset \mathbb{R} \) be a (bounded, closed) interval.
  Then
  \begin{equation}
  \label{eq:1}
  m^{*}([a,b]) \coloneq \inf \sum\limits_i l(I_i) = |b-a|,
  \end{equation}
  with the infimum being taken over all sequences \( (I_i)_i \) of open intervals such that \( \bigcup_i I_i \supset [a,b] \) and \( l(I) \) denoting the usual length of the interval \( I \).
\end{problem}
\begin{proof}
  First of all note that the \( \leq \) direction is straight forward, since we may extend \( [a,b] \), for any \(\varepsilon>0\), into an open interval \( (a - \frac{\varepsilon}{2}, b + \frac{\varepsilon}{2}) \).
  Then that cover has length \( |b-a| + \varepsilon \); \( m^{*} \) is a lower bound, so it must satisfy \( m^{*} \leq |b-a| + \varepsilon \) which, since \(\varepsilon\) was arbitrary, implies \( m^{*} \leq |b-a| \).

  In the other direction, let \( \{ I_j \}_1^{\infty} \) be an arbitrary open cover of \( [a,b] \).
  By the Heine-Borel property, \( [a,b] \) is compact; hence, there exists a finite subcover \( {I_1, \dots, I_N}\), \( N \in \mathbb{N} \).
  We can, without loss of generality, assume this cover is minimal, since we can simply remove a finite number of superfluous intervals if necessary.
  Then, writing \( I_j = (\alpha_j, \beta_j) \), we must have that
  \begin{equation}
  \label{eq:2}
  \alpha_1 < a < \beta_1, \quad \alpha_2 < \beta_1 < \beta_2, \quad \dots, \quad \alpha_N < \beta_{N-1} < b < \beta_N.
  \end{equation}
  The sum of the lengths is then
  \begin{equation}
  \label{eq:3}
  (\beta_1 - \alpha_1) + (\beta_2 - \alpha_2) + \dots + (\beta_N - \alpha_N) \geq \beta_N - \alpha_1 \geq b-a.
  \end{equation}
\end{proof}
\end{document}
